% Source-backed numbered records rebuilt from the visual page ledger.
% Statements and source proofs are kept in separate environments.

\begin{remark}[Source statement 7.1.5]\label{fga-rem-7-1-5}\dcref{7.1.5}\source{fga-explained:page-0172}\uses{}
\begin{verbatim}
Let X be a smooth connected varietyover k and let^ be a
coherentsheafon X with supp{T) ^ X. For an irreducible    hypersurfaceV (ZX
let \y\ be itsgenericpoint,and put R := Ox^[v]- Let M be the stalkof T at [V].
Then R isa,discretevaluationring,and by [Har77] IIL6.11A,1116.12A, M has
homologicaldimension 1, i.e.thereexistsa freeresolution
                            0 -> i?^ -^ i?^ -> M -> 0.
Note that the two modules on the lefthave the same rank, because supp{T) ^ X.
Let my G Z>o be the valuationof det((/?)G R.
    More geometricallywe can describethisas follows:There isan open subset
IJ<ZX whose intersection with V isopen, on which we have a resolution
                            0 _> (9^ -!^(9^ _> JT _> 0.
Then letmy be the order of vanishingof det{(p)on an open subsetof V.
    By the previousLemma, my isindependent of the choiceof the resolution.
\end{verbatim}
\end{remark}
\begin{proof}
\notready The source states this result without supplying a proof at this location.
\end{proof}


\begin{definition}[Source statement 7.1.6]\label{fga-def-7-1-6}\dcref{7.1.6}\source{fga-explained:page-0172}\uses{fga-rem-7-1-5}
\begin{verbatim}
Under the assumptions of Remark 7.1.5,let

                              div{T) := y^myV.
                                         V
                                Cartierdivisoron X. The sum isfinitebecause
This is by definitionan effective
my can only be nonzero ifV C supp{T).
    We want to constructsuch a divisordiv{J~)in a relativesituation.Before
                                                 we collecta few usefulresults.
statingand proving the main theorem to thiseffect,
7.1.THE SYMMETRIC POWER AND THE HILBERT-CHOW MORPHISM               163
\end{verbatim}
\end{definition}
\begin{proof}
\notready The source states this result without supplying a proof at this location.
\end{proof}


\begin{definition}[Source statement 7.1.9]\label{fga-def-7-1-9}\dcref{7.1.9}\source{fga-explained:page-0173}\uses{}
\begin{verbatim}
Let y be a scheme. A (notnecessarily   closed)pointp of y
has depth d ifdepth Oy^p isd. In particularithas depth 1 ifrup C Oy^p containsa
non zerodivisorand forevery non zerodivisort G rripevery element inthe maximal
idealof the quotientring Oy^p/{t)isa zerodivisor.
\end{verbatim}
\end{definition}
\begin{proof}
\notready The source states this result without supplying a proof at this location.
\end{proof}


\begin{remark}[Source statement 7.1.10]\label{fga-rem-7-1-10}\dcref{7.1.10}\source{fga-explained:page-0173}\uses{}
\begin{verbatim}
Ify isa smooth variety,then pointsof depth 1 are generic
pointsof prime divisors,   closedirreducible
                       i.e.,               hypersurfaces.
   The proof of the followingLemma   can be found in [EGAIV4], Prop. 21.1.8.
\end{verbatim}
\end{remark}
\begin{proof}
\notready The source states this result without supplying a proof at this location.
\end{proof}


\begin{lemma}[Source statement 7.1.12]\label{fga-lem-7-1-12}\dcref{7.1.12}\source{fga-explained:page-0174}\uses{}
\begin{verbatim}
A) Let f be a freecoherentsheafof rank r on a scheme Y.
Then detf := A^S isisomorphicto Oy'-,the isomorphism isunique up to a section

      B) Let
                         0 ?> fn ?^ ^n-1 ?>??.?> fo ?^ 0
be an exact sequence of locallyfreecoherentsheaveson a scheme Y. Then thereis
a canonicalisomorphism

                             (^{det SiY-'^ -^Oy.
\end{verbatim}
\end{lemma}
\begin{proof}
\begin{verbatim}
A) istrivial.B) willbe proven by induction.Do firsthe casen = 2.
    If0 ?>    ?^  ?>   ?^
           f 2 f 1 f 0 0 isa shortexact sequence of locallyfreesheaves,then
thereexistsa canonicalisomorphism
                          detf 2 ^ detSq ^ det ?i.
In fact,locallyon an open setchoose a basis/i,.??,/rs?^i?????^ro of fi such that
the fi are the image of a basisof ?2 and the gi map to a basisof Sq (both denoted
by the same letters). Then the isomorphism isgivenby
         (/lA ...A /^J (8)(^1A ...A Qr^)H^ /i A ...A /^2A ^1 A ...A Qr^.
Itiseasy to verifythat thisisindependent of the choiceof the gi and thus gluesto
a globalisomorphism.
    For the inductionstep,note that we can splitthe given exact sequenceup into
two exact sequencesof locallyfreesheaves
               0->f->fi   ->fo->0,     0->fn ->...-> ^2 ->f->0
and obtain by inductiona canonicalisomorphism (^^^Q (detfi)         '^Oy.     D
\end{verbatim}
\end{proof}


\begin{remark}[Source statement 7.2.2]\label{fga-rem-7-2-2}\dcref{7.2.2}\source{fga-explained:page-0176}\uses{}
\begin{verbatim}
If X is a smooth irreduciblesurface5, one can show that
P = F{Iz^{s)) is just the blowup of 5 x S^'^^along the universalfamily Zn{S)
            There itisalsoshown that P(/^E)) isthe incidencecorrespondence
[E11],[ES98].
                5[n,n+i]_ I{Z,W) G 5^^^X 5^^+^^IZ c
                                                    Wy
7.2.IRREDUCIBILITY AND NONSINGULARITY                    167

Here Z C W means that Z is a subscheme of W. One often uses 5^^'"^"^^^
                                                                    to
understand propertiesof the Hilbertschemes of pointsinductively.
    Let X be a nonsingularquasiprojective   varietyof dimension d. Let Xq C X^
be the dense open setof (pi,...,pn) with the pi distinct.Let Xq be itsimage in
X^'^\ which parametrizes effectivez ero cyclesX]^[pi]with the pi distinct.This is
alsoopen and dense. As Sn acts freelyon (X^)o we see that Xq is nonsingular
of dimension nd. Let Xq be the preimage in X^'^^.One checksthat at any point
of Xq the dimension of the tangentspace isnd and that p\^[n]isan isomorphism.
Thus X^'^^containsa nonsingularopen subsetwhich isisomorphicto an open subset
ofX(^). In the case that X isa curve or a surfaceone can use thisto show that
     isnonsingularand irreducible.
\end{verbatim}
\end{remark}
\begin{proof}
\notready The source states this result without supplying a proof at this location.
\end{proof}


\begin{remark}[Source statement 7.2.4]\label{fga-rem-7-2-4}\dcref{7.2.4}\source{fga-explained:page-0178}\uses{}
\begin{verbatim}
Note that in thisproof we do not show that the obstruction
space ExtQ^{Xz^Oz) vanishes;in factitusuallywillnot.
\end{verbatim}
\end{remark}
\begin{proof}
\notready The source states this result without supplying a proof at this location.
\end{proof}


\begin{remark}[Source statement 7.2.5]\label{fga-rem-7-2-5}\dcref{7.2.5}\source{fga-explained:page-0178}\uses{}
\begin{verbatim}
Let X be a nonsingularvariety.Then X^'^^isnonsingularfor
n <3.
\end{verbatim}
\end{remark}
\begin{proof}
\begin{verbatim}
Let d = dim{X). Itisenough to show that homoxi^z, Oz) < dn for
all[Z] gXH.  Obviously
                 }iomox{^z,Oz)=      0      Homc)xp(Jz,p,C>z,p).
                                   pEsupp{Z)
Thus we can reduce to the casethat supp{Z) isa pointp. Let m be the maximal
idealat p. Is iseasy to show that there are localparameters xi,...,Xd at p such
that Oz,p isof the form

                    k[xi,...,Xd]/m                     n = 1,
                    k[xi,...,Xd]/{m'^+ {X2,...,Xd))    n = 2,


                   k[xi,...,Xd]/{m^+ {x2,...,Xd))    n = 3.
                                      + {xs,...,Xd))
                  [or k[xi,...,Xd]/{m'^
                                              ^z,p) = dlen{Oz,p)-
In allcasesone easilychecksthat Hottiq^ {^z.p-,                              [proof end]
\end{verbatim}
\end{proof}


\begin{remark}[Source statement 7.2.6]\label{fga-rem-7-2-6}\dcref{7.2.6}\source{fga-explained:page-0178}\uses{}
\begin{verbatim}
Let X be nonsingularof dimension 3. Then X^^^ issingular.
\end{verbatim}
\end{remark}
\begin{proof}
\begin{verbatim}
Let [Z] e X^^^ be the point Oz = Op/m?, Then
                Yioiaox{^z,Oz) = Hom/c (mVm^m/m^)     = k^^.
So the dimension isbiggerthan dn = 12.                                       D
\end{verbatim}
\end{proof}


\begin{remark}[Source statement 7.2.7]\label{fga-rem-7-2-7}\dcref{7.2.7}\source{fga-explained:page-0178}\uses{}
\begin{verbatim}
Let X be a nonsingularvariety.Note that a 0-dimensional
subscheme Z C X oi length at most 3 has embedding dimension 2, i.e.islocallya
subscheme of a nonsingularsurface.In factitisnot difficult
                                                        to seethat ifZ c Xt^]
has embedding dimension at most 2, then [Z] isa nonsingularpoint of X.
\end{verbatim}
\end{remark}
\begin{proof}
\notready The source states this result without supplying a proof at this location.
\end{proof}


\begin{remark}[Source statement 7.2.9]\label{fga-rem-7-2-9}\dcref{7.2.9}\source{fga-explained:page-0179}\uses{}
\begin{verbatim}
One can show that (see[Iar85])
     A) (P^)N isirreducible  forn < 7 and alld,
     B) (P^)[^]isirreducible forti= 3 and reducibleforti> 4.
                 to show that theseresultsalsohold forX any smooth quasipro-
It isnot difficult
jectivevarietyof dimension d insteadof P*^.

                        7.3. Examples   of Hilbert schemes
    Let X be a nonsingularprojectivevarietyof dimension d. We givesome
    of X^'^^forsmallvaluesof n.
\end{verbatim}
\end{remark}
\begin{proof}
\notready The source states this result without supplying a proof at this location.
\end{proof}


\begin{proposition}[Source statement 7.3.3]\label{fga-pro-7-3-3}\dcref{7.3.3}\source{fga-explained:page-0179}\uses{}
\begin{verbatim}
Let C be a nonsingularquasiprojective
CW    ?> G(^) isan isomorphism.
\end{verbatim}
\end{proposition}
\begin{proof}
\begin{verbatim}
As the localring of C at a point p is a discretevaluationring,all
idealsin Oc,p are powers of the maximal idealrrip.Thus forall[Z] G C^^^we have

                      oz = 0oc,p./<%                 = '
                                               E"^
170                       7.HILBERT SCHEMES OF POINTS

Then p sends Z to Xl^^i[Pi].Thus p isbijective.                    itisan
                                              As itisalsobirational,
isomorphism by Zariski'sM ain Theorem.
    Alternativelyone can seethat
               TT:C X C("-i) -^ C("\              ^
                                                    J^n^lpi] + [p]
                                    (p,^ni[pi])
isflatof degree n over C(^\ definingan inverseto p.                          D

    Now let5 be a nonsingularsurface.We have seen that the open subset 5o
of 5^^^ consistingof sums of n distinctpointsisnonsingular.Now we want to see
that 5^^^ \ S^^^ isthe singularlocus of 5^^). By Example 7.1.3thisistrue for
n = 2, because              isthe image of the diagonal in 5^ under the quotient
map  5^  ?> 5^^\  Let A  = Ui<i<7<n^^j ^^ ^^^ ^^S diagonalin 5^, where A^j =
{(j;i,...,j;^)^ 'S'^I ^i = ^j}? ^^d let Aq be the open subset where precisely2
of the Xi coincide.Let p G Aq; we can assume that p = (z,z,2:3,...,   z^) G Ai^2-
Then the stabilizer  of p in 5^ is {1,t} where r isthe transpositionof the first
and second factor.Thus by Example 7.1.3we see that a formal neighborhood of
the image of p in 5^^^ (i.e.the completed localring of 5^^^ at p) isisomorphicto
A:[[it,v^!(;,   j;3,ys,...,Xn^yr^jiuw? v'^).Thereforeallthe pointsin the image
          j;',y',
of Aq are singularpoints of 5^^^. The closureof the image of Aq is5^^^ \ 5o ?
As the singularlocusof 5^^^ is closed,itfollowsthat the singularlocusof 5^^^ is
precisely5^^) \ 5^""^
\end{verbatim}
\end{proof}


\begin{theorem}[Source statement 7.3.4]\label{fga-the-7-3-4}\dcref{7.3.4}\source{fga-explained:page-0180}\uses{}
\begin{verbatim}
(Fogarty [Fog68]) Let 5 be a nonsingularquasiprojective
surface.Then p :5^^^?> 5^^^ isa resolutionof singularities.
\end{verbatim}
\end{theorem}
\begin{proof}
\begin{verbatim}
5^ isnonsingularand irreducible,  and p isan isomorphism over the
open subset Sq , which isthe nonsingularlocusof 5^^^.Thus itisa resolutionof
singularities.                                                            D
\end{verbatim}
\end{proof}


\begin{remark}[Source statement 7.3.5]\label{fga-rem-7-3-5}\dcref{7.3.5}\source{fga-explained:page-0180}\uses{}
\begin{verbatim}
A) One of the reasonsforthe interestin the Hilbertschemes of pointson a
         surfaceisthat they give canonicalresolutionsof the singularities
                                                                        of the
         symmetric powers.
      B) In factitturns out [Bea83] that p : 5^^^ ?> 5^^^ has a particularnice
         property:5^^^ isGorensteinand Kg[n] = p''{Ks(n)). One says that p isa
         crepant resolutionof singularities.
      C) In case5 isa K3-surfaceor an Abelian varietyitisalsoshown in [Bea83]
         that S^'^^
                  ishomomorphic symplectic,i.e.itcarriesan everywhere nonde-
         generateholomorphic two-form. This isan additionalreason forinterest
         in Hilbertschemes of pointsbecause holomorphic symplecticvarietiesare
         very rare. In factalmost allknown holomorphic symplecticprojective
         varietiescan be relatedto Hilbertschemes of points.

                 7.4. A stratificationof the Hilbert schemes
    For the restof Chapter 7 let5 be a smooth projectivesurfaceoverthe complex
numbers.
    We want to study a naturalstratification of 5^^^ and 5^^^.
7.4.A STRATIFICATION OF THE HILBERT SCHEMES                 171
\end{verbatim}
\end{remark}
\begin{proof}
\notready The source states this result without supplying a proof at this location.
\end{proof}


\begin{definition}[Source statement 7.4.3]\label{fga-def-7-4-3}\dcref{7.4.3}\source{fga-explained:page-0181}\uses{}
\begin{verbatim}
Let Hn := Hilb'^{Spec{k[[x,y]]/{x,y)'^),
                                                           parametrizingthe
idealsof colength n in A:[[j;,y]]
                              (note that every idealof colength n in A:[[j;,y]]
contains{x^y)^).The Hn are calledthe punctual Hilbertschemes.
\end{verbatim}
\end{definition}
\begin{proof}
\notready The source states this result without supplying a proof at this location.
\end{proof}


\begin{remark}[Source statement 7.4.5]\label{fga-rem-7-4-5}\dcref{7.4.5}\source{fga-explained:page-0182}\uses{}
\begin{verbatim}
H^ isa locallytrivial
            of dimension n ? 1.
isirreducible
\end{verbatim}
\end{remark}
\begin{proof}
\begin{verbatim}
We have a morphism H^ -^ P-*^,     sending Z to itstangent space at
@,0). Let U cH^he    the inverseimage of  A^ = P^ \ (j;= 0). Then Z eU lieson
a curve C which is given as the graph of a power seriesy = Yli>i (^i^^-Thus we
can write
                     Iz = {y - dix - '??- an-ix''~^,x'').
Thus U is isomorphic to A^~-^ with coordinatesai,...,a^-i, and the map to A^
is given by (ai,...,a^-i) f-^ai. Similarlyone argues over the inverseimage of
pi \ (y = 0).                                                              D


    In the case n = 3 we see that the only subscheme which isnot curvilinearis
the scheme with ideal(x,y)^ and H? isdense in H3. In factthisistrue in general.

   Theorem                                 and H^ isopen and dense in Hn-
              7.4.6.[Bri77] Hn isirreducible

    Another proof of thisirreducibility
                                     in [Ell],[ES98] makes use of the factthat
one has an inductivedescriptionof the S^'^^
                                          via the incidencevarieties        =
                                                                    S^'^'^'^^'^
{{Z,W)  G 5f^^X 5f^+i^           = P{Izr.{S))(seeRemark  7.2.2).
                        \Z CW}
7.5.THE BETTI NUMBERS OF THE HILBERT SCHEMES OF POINTS            173

        7.5. The Betti numbers     of the Hilbert schemes of points
    Now I want to summarize what we have shown so farand put itintocontext.
Then we willuse itto compute the Bettinumbers of the Hilbertschemes of points
on a surface.In thissectionwe assume that 5 isa smooth projectivesurfaceover
C.
\end{verbatim}
\end{proof}


\begin{proposition}[Source statement 7.5.2]\label{fga-pro-7-5-2}\dcref{7.5.2}\source{fga-explained:page-0183}\uses{}
\begin{verbatim}
p : S^'^^-^ 5^"")is strictlysemismallwith respectto the
stratification
            by the Su - Furthermore the fibres of p are irreducible.
\end{verbatim}
\end{proposition}
\begin{proof}
\begin{verbatim}
We know that S^'^^is nonsingularand irreducibleof dimension 2n,
which isalsothe dimension of S^'^\For u = (ni,...,rir)a partitionof n, we have
seenthat Su isnonsingularofdimension 2r,i.e.itscodimension isJ2^i=i2(n^? 1).
We have shown that SlT -^ Sl^^ isa locallytrivialfibrebundle in the analytic
                                                           ofdimension rii? l.
topologywith fibreiJ^ix ...x Hn^ and the Hn, are irreducible
This shows the claim.                                                       D

    This resultcan be used to compute the Bettinumbers of the Hilbertschemes.
One can use the propertiesof intersection cohomology and perversesheavesto get
a short argument, almost entirelywithout computations. The disadvantageof this
approach isthat one has to use advanced toolsand results.We willhere describe
these tools very briefly.One can then use them as a black box, i.e. one can
determinethe Bettinumbers of the Hilbertschemes quiteeasilywithout having to
reallyunderstand what intersection  cohomology and perversesheavesare.
    Let Y be an algebraicvarietyover C We willonly use the analytictopology
over C and allcohomology considerediswith Q coefficients.   We write W{Y) :=
H\YM)-
    Let ICy be the intersection  cohomology complex on Y. This is a complex
of sheaveson Y (strictly  speaking an element in the derived category),with the
property thatitscohomology IWiY) := H'^{Y,ICy) isthe intersection    cohomology
of Y. The intersection  cohomology groups IWiY) are definedfor any algebraic
varietyY and fulfill
                   Poincaredualitybetween IH~'^{Y) and IH'^{Y).IfY issmooth
and projectiveof complex dimension n, then ICy = QyM        is just the constant
sheaf Q put in degree n. Thus in thiscase IH^~^{Y) = H^{Y). More generally
the same istrue ify = X/G isthe quotientof a smooth projectivevarietyby a
finitegroup. One of the most powerful resultsabout intersection cohomology and
perverse sheavesisthe decomposition  theorem of [BBD82] forpropermorphisms of
174                            7.HILBERT SCHEMES OF POINTS

algebraicvarieties. This resultbecomes much simplerforsemismallmorphisms and
in thiscase computes the intersection cohomology ofX interms ofthe intersection
cohomologies  of the closuresFq; of the strataof Y.
\end{verbatim}
\end{proof}


\begin{lemma}[Source statement 7.5.3]\label{fga-lem-7-5-3}\dcref{7.5.3}\source{fga-explained:page-0184}\uses{}
\begin{verbatim}
[GS93] Let / :X -^ y be a proper morphism, which isstrictly
semismall with irreducible                                   Y = JJ^,Yqc into
                         fibreswith respectto a stratification
locallyclosednonsingularstrata.Then  Rf^{ICx)  = Yla ^^y  ?
    Here Rf^ isthe pushforward in the derivedcategory,and Y^ isthe closureof
\end{verbatim}
\end{lemma}
\begin{proof}
\notready The source states this result without supplying a proof at this location.
\end{proof}


\begin{lemma}[Source statement 7.5.4]\label{fga-lem-7-5-4}\dcref{7.5.4}\source{fga-explained:page-0184}\uses{}
\begin{verbatim}
[GS93] Let n : X -^ Y he a,finitebirationalmorphism of
irreducible                 Then Rt:^{ICx) = ICy-
          algebraicvarieties.
    Now we want to use these two resultsto compute the Betti numbers of the
Hilbertschemes. Write bi{Y) := dimH'^{Y) forthe i-thBetti number of Y and
p{Y) := J2i^i{Y)z'^forthe Poincare polynomial.
\end{verbatim}
\end{lemma}
\begin{proof}
\notready The source states this result without supplying a proof at this location.
\end{proof}
